\ifdefined\pdfinfoomitdate
  \pdfinfoomitdate=1
\fi
\ifdefined\pdfsuppressptexinfo
  \pdfsuppressptexinfo=15
\fi
\documentclass[10pt]{article}
\usepackage[letterpaper,margin=1in]{geometry}
\usepackage[T1]{fontenc}
\usepackage[utf8]{inputenc}
\usepackage{lmodern}
\usepackage{microtype}
\usepackage{amsmath,amssymb}
\usepackage{xurl}
\usepackage[hidelinks]{hyperref}
\hypersetup{
  pdftitle={Supplementary Note: Positive Recurrence for Single-Linkage Bimolecular Weakly Reversible Stochastic Reaction Networks},
  pdfauthor={Alec Kriebel},
  pdfsubject={Trace-chain, physical-time, and computational verification interfaces},
  pdfkeywords={stochastic reaction networks, positive recurrence, supplementary note, reproducibility}
}
\setlength{\parindent}{1.2em}
\setlength{\parskip}{0pt}
\emergencystretch=2em

\title{\bfseries Supplementary Note:\\[0.25em]
Positive Recurrence for Single-Linkage Bimolecular\\
Weakly Reversible Stochastic Reaction Networks}
\author{Alec Kriebel\\[0.2em]
\small Independent researcher\\
\small ORCID: \href{https://orcid.org/0009-0001-9320-500X}{0009-0001-9320-500X}}
\date{Supplementary Note --- 16 August 2026}

\begin{document}
\maketitle

\begin{abstract}
This supplementary note collects edge-case clarifications and reproducibility
interfaces supporting the main manuscript.  The finite-trace and
embedded-chain-to-continuous-time propositions, including their proofs, remain
in the main manuscript.
\end{abstract}

\section{Trace-chain and physical-time details}

\subsection{Labelled-channel marking}

The augmented state records the target of the actual labelled channel that
fired.  If two channels have the same population displacement but different
sources or targets, their post-jump populations can agree while their carried
targets or transition probabilities differ.  They therefore remain separate.
Only channels with identical source and target may be aggregated.

\subsection{Nonemptiness of the Foster set}

The finite set $K$ in the episode-drift construction is nonempty, not merely
finite.  The global-minimum argument is pathwise: the residual log-factorial
potential has a global minimum, and every possible episode endpoint has
potential at least that minimum.  No optional-stopping or integrability claim
is needed at this step because every episode has uniformly bounded finite
length.

\subsection{Hitting versus positive return}

The trace-chain proof distinguishes first hitting from positive return.  From
$k\in K$, one ordinary jump is taken before the next hit of $K$, so the
resulting return time is positive.  There are only finitely many successors of
states in $K$, and the endpoint-chain drift estimate gives finite expected
hitting of $K$ from each successor.  The finite trace chain is considered only
after this step.  A uniform finite-block chance of reaching one selected trace
state gives finite expected positive trace return; the random trace excursions
are then expanded into ordinary embedded jumps using the strong Markov
property and Tonelli's theorem.

\subsection{Stationary occupation balance}

The return construction gives a finite expected number of embedded jumps and
a finite expected amount of physical time before the positive return to the
selected state $x_*$.  During a complete return cycle, the number of arrivals
at any state $y\ne x_*$ equals the number of departures from $y$ pathwise.  At
$x_*$ the initial departure and terminal arrival also balance.  Taking
expectations and weighting each visit by its exponential holding-time mean
gives the statewise stationary balance equations for the normalized mean
occupation measure displayed in the manuscript.  This is the elementary
cycle-level interface behind the cited regenerative occupation theorem.

\subsection{Nonexplosion in physical time}

The population embedded chain is constructed before nonexplosion is known.
The finite-mean return result for that embedded chain supplies a population
state $x_*$ visited infinitely often.  At each visit, the following holding time has the same exponential law
with one fixed finite rate.  The strong Markov construction makes these
holding times independent copies, and their infinite sum diverges almost
surely.  Total physical time dominates this subseries, so jump times cannot
accumulate at a finite time.  This recovers nonexplosion for the present
subclass; broader bimolecular weakly reversible nonexplosion is already known,
as cited in the main manuscript.

\subsection{Absorbing singletons}

Absorbing singleton reachability classes are handled before channel
augmentation and carry their point-mass stationary laws.  In a nonabsorbing
irreducible class every population state enables a genuine channel; otherwise
that state would itself be absorbing, and symmetry of accessibility would
force its reachability class to be the singleton.

\section{Computational verification boundary}

The accompanying dependency-free Python package performs deterministic,
exact-arithmetic checks of the falling-factorial identity and
source-probability rewrite.  It also tests state-cycle lifting and accessibility
symmetry on finite calibration networks, scalar-envelope branch conditions and
monotonicity, the bimolecular top-complex classification with
certificate-validated witnesses, absorbing-singleton handling, the corrected
$\kappa_2\downarrow0$ rate-example limit, the exact logarithmic coefficient,
the exact Anderson--Cappelletti--Kim Example~4.1 comparison included in the
paper, and the stationary return-cycle normalization on a finite chain.

The reproducer runs the complete test suite, emits the canonical report twice,
requires those reports to be byte-identical, and compares the result with the
committed report.  Environment information is recorded separately so that it
does not destabilize the canonical JSON.  The canonical report is verified on
CPython 3.11 through 3.14; no claim is made for future interpreter versions
until they are added to the test matrix.

These computations are falsification and regression aids.  No finite atlas
proves the universal theorem, no random test proves recurrence, and no
computation enumerates the analytic Foster set $K$ or certifies a useful bound
on its cardinality, location, or diameter.  Those universal conclusions come
only from the manuscript's analytic argument.

\section{Reproduction pointers}

From the tagged Version~1.2.1 release directory, run:

\begin{verbatim}
cd code
./reproduce.sh
\end{verbatim}

The canonical output is
\path{code/verification_report.json}; an identical durable copy is retained
as \path{supplement/verification_report.json}.  The reproducibility environment,
release identifiers, and complete clean-checkout replay procedure are recorded
in \path{REPRODUCIBILITY.env} and \path{validation/}; the release manifest is
\path{supplement/MANIFEST.sha256}.  The package archive is created
deterministically by \path{supplement/build_release_archive.py} and checked
byte-for-byte in continuous integration.

\end{document}
